% Complementary outputs with dead time, drawn to scale, and the encoding that
% turns a requested time into a register value.
\begin{tikzpicture}[font=\sffamily\scriptsize, x=1cm, y=1cm]

% ================================================================ band 1: the pair
% The two dead-time windows are painted first, so the traces stay on top of them.
\foreach \a/\b in {1.8/2.4, 6.0/6.6}{ \fill[red!14] (\a,0.9) rectangle (\b,3.5); }

\node[signame] at (-0.3,3.0) {high side};
\node[signame] at (-0.3,1.3) {low side};

% high side: off, then on from 2.0 to 6.0, then off
\draw[signal, line width=1pt] (0,2.8) -- (2.4,2.8) -- (2.4,3.4) -- (6.0,3.4)
                              -- (6.0,2.8) -- (8.4,2.8);
% low side: on, then off from 1.8 to 6.2, then on again
\draw[signalb, line width=1pt] (0,1.7) -- (1.8,1.7) -- (1.8,1.1) -- (6.6,1.1)
                               -- (6.6,1.7) -- (8.4,1.7);
\node[tlabel, anchor=east] at (-0.35,3.4) {on};
\node[tlabel, anchor=east] at (-0.35,2.8) {off};
\node[tlabel, anchor=east] at (-0.35,1.7) {on};
\node[tlabel, anchor=east] at (-0.35,1.1) {off};

\foreach \a/\b in {1.8/2.4, 6.0/6.6}{ \draw[faultedge, <->] (\a,3.72) -- (\b,3.72); }
\node[tlabel, anchor=south, text=red!60!black] at (2.1,3.8) {dead time};
\node[tlabel, anchor=south, text=red!60!black] at (6.3,3.8) {dead time};
\node[tlabel, anchor=west, text width=34mm] at (8.6,3.0)
  {In the two shaded windows neither transistor is on. A bridge without them is
   a short across its own supply.};

% ================================================================ band 2: measured how
\node[chlabel, anchor=west] at (0,0.35) {how it gets measured here, with no oscilloscope};
\draw[busstub] (2.4,2.7) -- (2.4,-0.1);
\draw[busstub] (1.8,1.0) -- (1.8,-0.1);
\draw[tspan] (1.8,-0.1) -- (2.4,-0.1);
\node[tlabel, anchor=west, text width=74mm] at (2.8,-0.1)
  {two capture channels and two jumper wires: the falling edge of one output and
   the rising edge of the other, counted in timer ticks of 3.57 ns};

% ================================================================ band 3: the encoding
\node[chlabel, anchor=west] at (0,-1.0) {the register value is not a time: the encoding is piecewise};
\foreach \i/\r/\s in {0/{range 1}/{finest step},
                      1/{range 2}/{twice the step},
                      2/{range 3}/{eight times},
                      3/{range 4}/{sixteen times}}{
  \node[regcell, minimum width=30mm, text width=28mm, anchor=south west,
        fill=kercol] at (\i*3.1,-1.5) {\r\\{\tiny \s}};
}
\node[chlabel, anchor=north west, text width=58mm] at (0,-2.4)
  {Four ranges, each with a different resolution, and the exact boundaries are in
   the reference manual for this part. The firmware asks for a time in
   nanoseconds, encodes it, decodes the encoding back, and prints both, so the
   value it got is visible rather than assumed. That arithmetic is tested on a
   host over the whole range including every boundary.};

\node[umlnote, text width=54mm, anchor=north west] at (6.2,-2.4)
  {Asking for five hundred nanoseconds and being told five hundred is a
   coincidence of where the ranges fall. On this part the encoded value for that
   request means a few nanoseconds more, and the sweep in step 5 shows the same
   three nanosecond offset at every setting, which is the propagation difference
   between the two pins and the capture path rather than an error in the dead
   time. It is reported and not subtracted.};
\end{tikzpicture}
